\section*{第 \ref{ch:g3:first-electric-circuit} 章——第一个电路}

\begin{solution}{exo:g3:first-electric-circuit:1}
\omterm{def:g3:first-electric-circuit:battery}{电池}的\omterm{def:g3:first-electric-circuit:battery}{接线柱}是 $+$（正）和 $-$（负）。
\omterm{def:g3:first-electric-circuit:bulb}{灯泡}的接触点：灯脚底部的尖端，
以及灯脚的金属侧面。
\end{solution}

\begin{solution}{exo:g3:first-electric-circuit:2}
闭合：一条从\omterm{def:g3:first-electric-circuit:battery}{接线柱}到\omterm{def:g3:first-electric-circuit:battery}{接线柱}、没有断开的回路。断开：一条
中间某处有缺口的回路。\omterm{def:g3:first-electric-circuit:bulb}{灯泡}只在闭合的回路中发光。
\end{solution}

\begin{solution}{exo:g3:first-electric-circuit:3}
回程。一根导线只是半个回路：电流必须
从一个\omterm{def:g3:first-electric-circuit:battery}{接线柱}出发，经过\omterm{def:g3:first-electric-circuit:bulb}{灯泡}，\emph{回到另一个
\omterm{def:g3:first-electric-circuit:battery}{接线柱}}——丽亚需要第二根导线，从\omterm{def:g3:first-electric-circuit:bulb}{灯泡}的侧面接触点
连到 $-$ \omterm{def:g3:first-electric-circuit:battery}{接线柱}。
\end{solution}

\begin{solution}{exo:g3:first-electric-circuit:4}
\omterm{ex:g3:first-electric-circuit:switch}{开关}就像回路中的一扇门：闭合时接通\omterm{def:g3:first-electric-circuit:circuit}{电路}；
断开时造成缺口。闭合的\omterm{ex:g3:first-electric-circuit:switch}{开关}内部有一小块金属片，
把两根导线连在一起。
\end{solution}

\begin{solution}{exo:g3:first-electric-circuit:5}
由直线连接而成的回路：两根长短不等的短线（\omterm{def:g3:first-electric-circuit:battery}{电池}，
较长的一根为 $+$），一个闭合放平的吊桥符号（\omterm{ex:g3:first-electric-circuit:switch}{开关}），
以及一个带十字的圆圈（灯）。
\end{solution}

\begin{solution}{exo:g3:first-electric-circuit:6}
\omterm{def:g3:first-electric-circuit:battery}{电池}用旧了；连接松动（导线或\omterm{def:g3:first-electric-circuit:bulb}{灯泡}没有拧紧，
弹簧弯曲）；\omterm{def:g3:first-electric-circuit:bulb}{灯泡}损坏——它发光的灯丝断了，
使\omterm{def:g3:first-electric-circuit:bulb}{灯泡}内部出现缺口。
\end{solution}

\begin{solution}{exo:g3:first-electric-circuit:7}
按钮\emph{就是}\omterm{ex:g3:first-electric-circuit:switch}{开关}。按下时，金属片接通，
回路闭合，铃响了；松开时，金属片抬起，
缺口出现，铃声停止。
\end{solution}

\begin{solution}{exo:g3:first-electric-circuit:8}
\omterm{def:g3:first-electric-circuit:battery}{电池}的推动力很小，很安全；墙上插座的推动力
要大数百倍，可能伤人甚至致命。插座不能用来
做实验，任何时候都不行。
\end{solution}

\begin{solution}{exo:g3:first-electric-circuit:9}
会亮。沿着回路走：$+$ \omterm{def:g3:first-electric-circuit:battery}{接线柱}，进入尖端，
穿过发光的灯丝，从侧面接触点出来，沿导线，
回到 $-$。这是一条完整的往返路径，
经过两个\omterm{def:g3:first-electric-circuit:battery}{接线柱}和两个接触点。
\end{solution}

\begin{solution}{exo:g3:first-electric-circuit:10}
他的回路从未经过第二个\omterm{def:g3:first-electric-circuit:battery}{接线柱}。\omterm{def:g3:first-electric-circuit:circuit}{电路}必须从
一个\omterm{def:g3:first-electric-circuit:battery}{接线柱}\emph{经过\omterm{def:g3:first-electric-circuit:bulb}{灯泡}到达另一个}\omterm{def:g3:first-electric-circuit:battery}{接线柱}：从
$-$ 出发又回到 $-$，是一条通往无处的往返路径——
\omterm{def:g3:first-electric-circuit:battery}{电池}的推动力需要它两个不同的端点都参与。
\end{solution}

\begin{solution}{exo:g3:first-electric-circuit:11}
所有\omterm{def:g3:first-electric-circuit:bulb}{灯泡}必须接在同一个回路上，一个接一个：
缺少一个\omterm{def:g3:first-electric-circuit:bulb}{灯泡}就会留下缺口，一处缺口就会使整串灯都熄灭——
这就是回路定律。（明年，这种一个接一个的连接方式
将有一个名称：串联\omterm{def:g3:first-electric-circuit:circuit}{电路}。）
\end{solution}
